\documentclass[11pt]{amsart}

\usepackage[utf8]{inputenc}

\usepackage{amssymb,amsmath,amsthm,mathtools}
\usepackage{enumitem}
\usepackage{booktabs,longtable,array}
\usepackage[margin=1.15in]{geometry}
\usepackage{xcolor}
\usepackage{url}
\usepackage[colorlinks=true,linkcolor=blue!50!black,citecolor=green!35!black,urlcolor=blue!50!black]{hyperref}
\usepackage[capitalise,nameinlink]{cleveref}

% ---------------------------------------------------------------------------
% Repository location. Every Lean name in the text refers to this tagged tree.
\newcommand{\repourl}{https://github.com/ndp75/dinfty}
\newcommand{\repotag}{v1.0.0}
\newcommand{\repodoi}{10.5281/zenodo.23423880}
% ---------------------------------------------------------------------------

% Lean declaration names: typewriter, breakable at dots.
\DeclareUrlCommand\lean{\urlstyle{tt}}
\makeatletter\g@addto@macro\UrlBreaks{\do\_}\makeatother
\emergencystretch=3em

\theoremstyle{plain}
\newtheorem{theorem}{Theorem}[section]
\newtheorem{proposition}[theorem]{Proposition}
\newtheorem{lemma}[theorem]{Lemma}
\newtheorem{corollary}[theorem]{Corollary}
\newtheorem*{theoremA}{Theorem A}
\newtheorem*{theoremB}{Theorem B}
\newtheorem*{theoremS}{Theorem}
\theoremstyle{definition}
\newtheorem{definition}[theorem]{Definition}
\newtheorem{remark}[theorem]{Remark}
\newtheorem{example}[theorem]{Example}

% Notation
\newcommand{\End}{\operatorname{End}}
\newcommand{\SEnd}{\operatorname{End}_\sigma}
\newcommand{\Glue}{\operatorname{Glue}}
\newcommand{\Aut}{\operatorname{Aut}}
\newcommand{\Fix}{\operatorname{Fix}}
\newcommand{\J}{\mathcal{J}}
\newcommand{\cst}[1]{c_{#1}}
\newcommand{\tg}[2]{\langle #1\rangle_{\mathsf{#2}}}
\newcommand{\Dc}{D_\infty}
\newcommand{\Db}{\overline{D}_\infty}
\newcommand{\BB}{\mathcal{B}}
\newcommand{\GG}{\Gamma}
\newcommand{\id}{\mathrm{id}}
\newcommand{\iso}{\cong}

\title[Rigidity and self-similarity in a gluing tower]{Rigidity and self-similarity in a three-copy gluing tower of finite lattices, machine-checked in Lean~4}

\author{Noah Pinto}
\address{Independent researcher}
\email{ndpinto1@gmail.com}

\subjclass[2020]{Primary 06B35; Secondary 06A06, 28A80, 68V20}
\keywords{domain theory, bilimit, rigidity, order automorphism, self-similarity, Sierpi\'nski gasket, iterated function system, Lean, formalized mathematics}

\begin{document}

\begin{abstract}
We study a tower of finite lattices $(D_k)_{k\ge 0}$ built by an order-theoretic
analogue of the three-copy gluing that presents the Sierpi\'nski gasket in
categories of metric spaces: $D_{k+1}$ is obtained by gluing three copies of the
lattice $\End(D_k)$ of monotone self-maps of $D_k$ at three seams. The levels
$D_k$ carry two towers of embedding--projection pairs, which give rise to three
limit objects: the colimit $\Dc$ of the constant embeddings, the bilimit $\Db$
of the same embeddings, and the bilimit $\BB$ of a second tower whose
embeddings are obtained by applying the gluing functor to the previous one.

We prove that $\Dc$ and $\Db$ are \emph{rigid}: each has the identity as its only
order automorphism (Theorem~B). We prove that $\BB$ is a fixed point of the
gluing built from its own Scott-continuous self-maps, $\Phi_\sigma(\BB)\iso\BB$
(Theorem~A). We show that $\Dc$ is exactly the poset of compact elements of
$\Db$, and that $\BB$ is order-isomorphic to neither $\Dc$ nor $\Db$, so the two
theorems are about genuinely different objects. On $\BB$ we build an internal
iterated function system with injective maps and relate its attractor to the
Sierpi\'nski gasket; here several natural expectations fail, and we prove the
failures: the comparison map onto the gasket is surjective and chain-continuous
but neither injective nor continuous for any topology with upward-closed open
sets, and the address coding map onto the gasket is a quotient map but not an
open map.

Every result is formalized in Lean~4 with Mathlib. The development has 116
modules and about 37{,}000 lines, and contains no \texttt{sorry}, no
\texttt{native\_decide} and no axioms beyond Lean's three standard ones.
\end{abstract}

\maketitle

\section{Introduction}\label{sec:intro}

Scott's construction of a non-trivial domain $D$ with $D\iso[D\to D]$
\cite{Scott1972} is the prototype for a broad family of results in which a
space is obtained as the limit of a tower of approximations and then shown to
be a fixed point of the functor that generates the tower
\cite{SmythPlotkin1982,AbramskyJung1994}. A different line of work presents
fractals as fixed points: Freyd showed that the unit interval is a final
coalgebra on bipointed sets \cite{Freyd2008}, Leinster gave a general theory of
self-similar spaces \cite{Leinster2011}, and Moss, Ratnayake and coauthors
presented the Sierpi\'nski gasket in categories of tripointed metric spaces by
completing the initial algebra of a three-copy gluing functor
$X\mapsto M\otimes X$ \cite{MossRatnayakeRose2014,RatnayakeEtAl2024}.

This paper combines the two. We take the gasket's three-copy, three-seam gluing
pattern and apply it not to a metric space but to the lattice of monotone
self-maps of a finite lattice. Iterating from the three-element chain gives a
tower of finite lattices $D_0, D_1, D_2,\dots$ that grows extremely fast, and
whose structure mixes Scott's function-space recursion with the gasket's
combinatorics. We ask two questions about its limits: are they
\emph{self-similar}, in the sense of being fixed points of the gluing; and are
they \emph{rigid}, in the sense of having no non-trivial order automorphisms?

\subsection*{Results}
The levels $D_k$ carry two towers of embedding--projection pairs (ep-pairs),
described in \cref{sec:objects}. They give three limit objects:
\begin{itemize}[leftmargin=2em]
\item $\Dc$, the colimit of the $D_k$ along the \emph{constant embeddings}
  $x\mapsto\tg{\cst x}{A}$;
\item $\Db$, the ep-pair bilimit of the same tower;
\item $\BB$, the ep-pair bilimit of the \emph{$\Phi$-iterate tower}, whose
  $(n{+}1)$-st embedding is the image under the gluing functor of the $n$-th.
\end{itemize}

\begin{theoremB}[\cref{thm:B}]
$\Aut(\Dc,\le)=\{\id\}$ and $\Aut(\Db,\le)=\{\id\}$. Moreover every level is
rigid, and every monotone injection $D_k\to D_k$ is the identity.
\end{theoremB}

\begin{theoremA}[\cref{thm:A}]
Let $\Phi_\sigma(\BB)$ be the three-copy gluing of the lattice of
Scott-continuous self-maps of $\BB$, at the constant map with value the
interior point of $\BB$. Then $\Phi_\sigma(\BB)\iso\BB$ as posets.
\end{theoremA}

\begin{theoremS}[\cref{thm:compacts,thm:gap}]
The canonical map $\Dc\to\Db$ is an order embedding onto exactly the compact
elements of $\Db$, and every element of $\Db$ is the directed supremum of the
compact elements below it; in particular $\Dc\not\iso\Db$. Further,
$\BB\not\iso\Db$ and $\BB\not\iso\Dc$.
\end{theoremS}

So Theorem~A and Theorem~B concern different objects, and neither transfers to
the other. Whether $\BB$ is rigid remains open; \cref{sec:Brigid} records the
reductions and refutations we do prove about it.

Finally, Theorem~A yields an iterated function system $(\sigma_A,\sigma_B,\sigma_C)$
of injective maps on $\BB$, and an address map whose image, the
\emph{attractor}, maps onto the geometric gasket $\GG$. The comparison is
instructive mainly for what fails. A first candidate system has the wrong
fixed-point set; a repaired one has the expected $\Fix(\sigma_A)=\{\bot\}$. The
map $\rho$ from the attractor onto $\GG$ is surjective, intertwining and
chain-continuous, but it is not injective, because the copies of $\BB$ meet at
no seam while those of $\GG$ do, and it is not continuous for any topology on
the attractor whose open sets are upward closed (in particular the Scott
topology). The standard address coding map onto $\GG$ is a quotient map but
not an open map. Each of these negative statements is a theorem
(\cref{sec:ifs}).

\subsection*{Formalization}
Every statement in this paper, apart from one clearly marked interpretive
remark, is a Lean~4 \cite{Lean4} declaration checked against Mathlib
\cite{Mathlib2020}. The release has 116 modules and 37{,}109 lines. A single
file, \lean{Dinfty/Verify.lean}, restates each headline result as an anonymous
\texttt{example} against the proved declaration, and prints the axioms it
depends on; all 195 reports are subsets of
$\{\texttt{propext},\texttt{Classical.choice},\texttt{Quot.sound}\}$.
Throughout, a Lean name in typewriter font, such as
\lean{CopyCHigh.colimit_compat}, refers to the tagged release at
\url{\repourl} (tag \texttt{\repotag}, archived at
\href{https://doi.org/\repodoi}{doi:\repodoi}); \cref{app:lean} collects them.
\Cref{sec:formal} describes the development and how to check it.

\subsection*{Contributions}
\begin{enumerate}[leftmargin=2em]
\item An order-theoretic gluing construction modelled on the gasket's functor
  $M\otimes-$, with $\End$ in place of the metric factor, and the tower of
  finite lattices it generates (\cref{sec:glue}).
\item Rigidity of the colimit and of the bilimit of that tower (Theorem~B).
  The colimit step is the main technical work; it replaces a level-detecting
  invariant, which cannot exist (\cref{prop:nofootholds}), by an exact
  combinatorial comparison of principal ideals.
\item A second ep-tower on the same levels whose bilimit is a fixed point of
  the Scott-continuous gluing (Theorem~A), together with proofs that it differs
  from the first tower's limits.
\item An internal IFS on that bilimit, and machine-checked negative results
  delimiting how far it can be identified with the geometric gasket.
\item A complete, axiom-audited Lean~4 formalization of all of the above.
\end{enumerate}

\subsection*{Organization}
\Cref{sec:related} discusses related work. \Cref{sec:glue} defines the gluing
and the tower. \Cref{sec:objects} defines the three limit objects and proves
they are pairwise distinct. \Cref{sec:B,sec:A} prove Theorems~B and~A.
\Cref{sec:ifs} treats the internal IFS and the gasket, and \cref{sec:Brigid}
what is known about automorphisms of $\BB$. \Cref{sec:formal} describes the
formalization.

\section{Related work}\label{sec:related}

\emph{Fixed points of domain constructors.} Scott's $D_\infty$
\cite{Scott1972} and its categorical reformulation via limit--colimit
coincidence \cite{SmythPlotkin1982} are standard; see
\cite{AbramskyJung1994,GierzEtAl2003}. Our Theorem~A follows the same pattern
for a mixed constructor: a function-space step followed by a three-copy gluing.
Two features differ from the classical setting. The gluing identifies points,
so its action on morphisms is only defined on maps preserving the three
distinguished points, and the tower has to be built so that its embeddings do.
And the function-space step must be taken over Scott-continuous maps for the
colimit argument to go through.

\emph{Fractals as fixed points.} Freyd's theorem \cite{Freyd2008}, Leinster's
general theory \cite{Leinster2011}, and Hasuo, Jacobs and Niqui's coalgebraic
representation theory of fractals \cite{HasuoJacobsNiqui2010} present
self-similar spaces as terminal coalgebras on categories of sets, topological
spaces or metric spaces. Moss, Ratnayake and Rose \cite{MossRatnayakeRose2014}
and Ratnayake, Manokaran, Jayewardene, Noquez and Moss \cite{RatnayakeEtAl2024}
obtain the gasket by Cauchy-completing the initial algebra of the gluing functor
$M\otimes-$ on tripointed metric spaces; Noquez and Moss treat the carpet
\cite{NoquezMoss2021}. We use their gluing pattern (three copies $a,b,c$; seams
$a.L\sim b.T$, $a.R\sim c.T$, $b.R\sim c.L$ in their notation) but replace the
metric factor by a function-space lattice, and ask order-theoretic rather than
metric questions.

\emph{Domain theory and IFS.} Edalat \cite{Edalat1996} represents IFS attractors
in power domains and upper spaces, following Hutchinson's construction
\cite{Hutchinson1981}. Our internal IFS on $\BB$ is in this spirit, but the
underlying domain is not a hyperspace of the plane; its relation to $\GG$ is
established through an explicit comparison map, and the negative results of
\cref{sec:ifs} show that this map is much weaker than an isomorphism.

\emph{Rigidity.} Rigidity of product-like structures is classically handled by
unique factorization: Sabidussi and Vizing's theorem for the Cartesian product of
connected graphs \cite{Sabidussi1960,Vizing1963}. The finite-level step of
Theorem~B has this shape, but we do not use it: the needed splitting follows from
counting atoms and cardinalities (\cref{lem:prodaut}).

\section{The gluing construction and the tower}\label{sec:glue}

All lattices are bounded. For a finite lattice $P$, write $\End(P)$ for the
lattice of monotone maps $P\to P$ under the pointwise order, and $\cst x$ for
the constant map with value $x$. When $P$ is a dcpo, $\SEnd(P)\subseteq\End(P)$
is the subposet of Scott-continuous maps (\lean{GluingTowerLibrary.ScottEndP}).

\begin{definition}[Three-copy gluing]\label{def:glue}
Let $Q$ be a finite lattice and $m\in Q$ with $\bot<m<\top$. Let $\Glue(Q,m)$ be
the set of three tagged copies $\tg{x}{A},\tg{x}{B},\tg{x}{C}$ ($x\in Q$),
modulo the three seam identifications
\[
  \tg{\bot}{B}=\tg{m}{A},\qquad \tg{\bot}{C}=\tg{\top}{A},\qquad \tg{m}{C}=\tg{\top}{B}.
\]
Define a relation $R$ on tagged points by
\begin{gather*}
  \tg xI\mathrel R\tg yI \iff x\le y\quad(I\in\{A,B,C\}),\qquad
  \tg xA\mathrel R\tg yB\iff x\le m,\\
  \tg xA\mathrel R\tg yC\ \text{always},\qquad
  \tg xB\mathrel R\tg yC\iff m\le y,
\end{gather*}
and no other instances. For $u,v\in\Glue(Q,m)$ put $u\le v$ if some
representative of $u$ is $R$-related to some representative of $v$.
\end{definition}

In the formalization the carrier is the set of canonical representatives
(\lean{GluingTowerLibrary.GlueCarGen}), each point has at most two
representatives, and the order is \lean{GluingTowerLibrary.leGlue}, with rule
\lean{GluingTowerLibrary.rawRule}. The library proves that $\Glue(Q,m)$ is a
finite lattice with bottom $\tg\bot A$ and top $\tg\top C$, and that each copy
inclusion $x\mapsto\tg xI$ is an order embedding. Copy $A$ sits at the bottom,
copy $C$ at the top, and copy $B$ is attached between them only through the
seams. The three distinguished points of $\Glue(Q,m)$ are
\[
  \bot=\tg\bot A,\qquad \tg mA,\qquad \top=\tg\top C .
\]

\begin{remark}
\Cref{def:glue} is the poset version of the gasket functor $M\otimes X$ of
\cite[Section~2.2]{RatnayakeEtAl2024}: three copies glued at three seams, each seam
identifying a distinguished point of one copy with a distinguished point of
another. What is new is the factor being glued (a lattice of self-maps rather
than a metric space) and the order (\cref{def:glue}) in place of a metric.
\end{remark}

\begin{definition}[The tower]\label{def:tower}
Let $D_0=\{\bot<\mu_0<\top\}$ be the three-element chain. Given a finite lattice
$D_k$ with an interior point $\mu_k$, set
\[
  D_{k+1}:=\Glue\bigl(\End(D_k),\,\cst{\mu_k}\bigr),\qquad
  \mu_{k+1}:=\tg{\cst{\mu_k}}{A}.
\]
\end{definition}

So the marked point at each level is the constant map at the previous interior
point, and the next interior point is that constant map seen in copy $A$
(\lean{SierpinskiApplication.markedPt}, \lean{SierpinskiApplication.dMid},
\lean{PhiTriTower.dMid_succ_eq}). The marked point satisfies
$\bot<\cst{\mu_k}<\top$ at every level
(\lean{SierpinskiApplication.D_marked_point_nondegenerate}), so the recursion is
well defined, and $\bigl|D_k\bigr|$ is strictly increasing in $k$
(\lean{TowerCardGrowth.card_D_strictMono}). The identity $D_{k+1}=
\Glue(\End(D_k),\cst{\mu_k})$ holds definitionally in the formalization
(\lean{PhiTriTower.D_succ_eq}).

\begin{remark}[Why a constant marked point]
On the two-element chain, $\End$ has exactly three elements and the identity
\emph{is} the middle one, so ``glue at $\id$'' and ``glue at a constant'' cannot
be told apart there. On the three-element chain they differ. We use the
constant throughout; it is the choice under which the gluing functor's action
on the tower's own maps preserves all three distinguished points
(\cref{sec:objects}).
\end{remark}

\begin{definition}[Tripointed maps; the action of $\Phi$]\label{def:phi}
A monotone map $f\colon(Q,m)\to(Q',m')$ is \emph{tripointed} if $f(\bot)=\bot$,
$f(m)=m'$ and $f(\top)=\top$. For such $f$ let $\Phi(f)\colon\Glue(Q,m)\to
\Glue(Q',m')$ act tagwise, $\tg xI\mapsto\tg{f(x)}I$.
\end{definition}

The tagwise action respects the seams precisely because $f$ is tripointed, and
$\Phi(f)$ is again tripointed. On tripointed maps $\Phi$ is a functor
(\lean{PhiTriFunctor.phiFun_id}, \lean{PhiTriFunctor.phiFun_comp},
\lean{PhiTriFunctor.phi_functor_laws}); on injective tripointed maps its action
coincides with the library's own gluing map
(\lean{PhiTriTower.phiFun_eq_glueMap}), with no further hypothesis.

\section{Three limit objects}\label{sec:objects}

\subsection{The constant tower: \texorpdfstring{$\Dc$}{D-infinity} and \texorpdfstring{$\Db$}{its bilimit}}

The \emph{constant embedding} is $\iota_k\colon D_k\to D_{k+1}$,
$\iota_k(x)=\tg{\cst x}A$ (\lean{SierpinskiApplication.iotaD}). It has an
explicit upper adjoint $p_k\colon D_{k+1}\to D_k$ with $p_k\iota_k=\id$ and
$\iota_kp_k\le\id$ (\lean{SierpinskiApplication.pD},
\lean{SierpinskiApplication.pD_iotaD},
\lean{SierpinskiApplication.pD_deflation}), so $(\iota_k,p_k)_k$ is a tower of
ep-pairs.

\begin{definition}
$\Dc$ is the colimit of $D_0\xrightarrow{\iota_0}D_1\xrightarrow{\iota_1}\cdots$
in posets: the quotient of $\coprod_kD_k$ identifying $x$ with $\iota_k(x)$
(\lean{SierpinskiApplication.DInfty}). $\Db$ is the limit along the projections:
sequences $(x_k)_k$ with $p_k(x_{k+1})=x_k$, ordered pointwise
(\lean{SierpinskiApplication.DBilimit}).
\end{definition}

The canonical map $\Dc\to\Db$ sends $x\in D_k$ to the sequence obtained by
projecting $x$ down to lower levels and embedding it into higher ones
(\lean{SierpinskiApplication.toDBilimit}).

\begin{theorem}\label{thm:compacts}
The map $\Dc\to\Db$ is an order embedding whose image is exactly the set of
compact elements of $\Db$. For every $y\in\Db$, the compact elements below $y$
form a directed set with supremum $y$. The top element of $\Db$, the sequence of
tops, is not in the image.
\end{theorem}

\begin{proof}[Lean]
\lean{SierpinskiApplication.toDBilimitEmb},
\lean{SierpinskiApplication.mem_range_toDBilimit_iff_isCompactElement},
\lean{SierpinskiApplication.DBilimit_compacts_le_directedOn},
\lean{SierpinskiApplication.DBilimit_isLUB_compacts_le},
\lean{SierpinskiApplication.topSeq_notMem_range_toDBilimit}.
\end{proof}

\begin{remark}
In classical terms $\Db$ is an algebraic domain whose basis of compact elements
is a copy of $\Dc$, i.e.\ the ideal completion of $\Dc$. We prove the four
facts of \cref{thm:compacts} but do not state the isomorphism
$\Db\iso\operatorname{Idl}(\Dc)$ as a separate theorem; this remark is the one
statement in the paper that is not itself a Lean declaration.
\end{remark}

\begin{proposition}[$\Dc$ has no footholds]\label{prop:nofootholds}
$\Dc$ has no atom and no maximal element: every $z$ has some $w>z$.
\end{proposition}
\begin{proof}[Lean]
\lean{DInftyNoFootholds.no_isAtom}, \lean{DInftyNoFootholds.no_top},
\lean{DInftyNoFootholds.no_footholds}.
\end{proof}

\Cref{prop:nofootholds} is what makes the colimit step of Theorem~B hard: any
invariant anchored at the bottom or at the top (heights, ranks, down-set
counts) works at each finite level and fails to detect levels in $\Dc$.

\subsection{The \texorpdfstring{$\Phi$}{Phi}-iterate tower and \texorpdfstring{$\BB$}{B}}

The constant embedding is not tripointed: $\iota_k(\top)=\tg{\cst\top}A\neq
\top$. So $\Phi$ cannot be applied to it, and the constant tower is not
generated by $\Phi$. We therefore build a second ep-tower on the levels
$D_1,D_2,\dots$ whose embeddings are tripointed.

\emph{Base step.} Define $e_0\colon D_1\to D_2$ tagwise on copies,
\[
  e_0\tg fA=\tg{\iota_0\circ f\circ p_0}A,\qquad
  e_0\tg fB=\tg{\lambda f}B,\qquad e_0\tg fC=\tg{\lambda f}C,
\]
where $\lambda f\in\End(D_1)$ is the explicit lift
\[
  (\lambda f)(w)=\iota_0\bigl(f(p_0w)\bigr)\ \vee\
  [\cst{\mu_0}\le f]\,\mu_1\ \vee\ [\top\le f]\,\top
\]
(Iverson brackets: the term is present if the condition holds and is $\bot$
otherwise). This is \lean{SierpinskiApplication.glueAA} at level $0$, with the
lift \lean{GlueAAJoin.liftT}. It is tripointed and preserves binary joins
(\lean{GlueAAJoin.glueAA_sup}); join preservation uses join-primality of the
marked point at every level (\lean{JoinPrimality.markedPt_joinPrime}), proved by
induction on $k$, since at level $k{+}1$ the interior point lies below every
point represented in copy $B$ or $C$. Its upper adjoint $q_0$ is constructed
explicitly, as a greatest element, and satisfies the
Galois condition $e_0(y)\le z\iff y\le q_0(z)$
(\lean{QAAExplicit.qAA_isGreatest}, \lean{QAAExplicit.glueAA_le_iff_le_qAA}).

\emph{Inductive step.} Given a tripointed ep-pair $(e_n,q_n)$ from $D_{n+1}$ to
$D_{n+2}$, conjugation $h\mapsto e_n\circ h\circ q_n$ and $h\mapsto q_n\circ h
\circ e_n$ are tripointed maps between the $\End$ lattices, and
\[
  e_{n+1}:=\Phi(h\mapsto e_nhq_n),\qquad q_{n+1}:=\Phi(h\mapsto q_nhe_n)
\]
is again a tripointed ep-pair, by functoriality of $\Phi$
(\lean{PhiIterateTower.phiStep}, \lean{PhiIterateTower.emb_succ}).

\begin{definition}
$\BB$ is the ep-pair bilimit of $(e_n,q_n)_n$: sequences $(x_n)_n$ with $x_n\in
D_{n+1}$ and $q_n(x_{n+1})=x_n$, ordered pointwise
(\lean{PhiIterateTower.phiTowerEP}, \lean{PhiBilimit.B}). It has three
distinguished points: $\bot$, the sequence $\mu_\BB$ of interior points, and the
sequence $\top_\BB$ of tops (\lean{PhiBilimit.midB}, \lean{PhiBilimit.topB}).
\end{definition}

\begin{remark}[How the towers meet]\label{rem:tripointed}
The constant tower is carried inside the $\Phi$-iterate tower:
$e_k\circ\iota_k=\iota_{k+1}\circ\iota_k$
(\lean{StageRigidity.legacy_subtower}). But no $\Phi$-iterate embedding of a top
is ever a constant-tower point: $e_n(\top)\neq\iota_{n+1}(x)$ for every $x$
(\lean{StageRigidity.phi_emb_top_not_legacy}). The two towers share their
finite stages and differ in how they glue them together.
\end{remark}

\subsection{The three objects are distinct}

Call $x$ in a poset \emph{$\omega$-compact} if whenever $c$ is a monotone
sequence with supremum $s\ge x$, already $x\le c_n$ for some $n$
(\lean{BilimitGap.OmegaCompact}). Order isomorphisms preserve
$\omega$-compactness.

\begin{theorem}\label{thm:gap}
$\BB\not\iso\Db$ and $\BB\not\iso\Dc$.
\end{theorem}

\begin{proof}
Every $e_n$ preserves $\top$, so a point of $\BB$ whose level-$0$ component is
the top of $D_1$ is $\top_\BB$; since the level-$0$ components of a chain in
$\BB$ eventually stabilize ($D_1$ is finite), $\top_\BB$ is $\omega$-compact
(\lean{BilimitGap.topB_omegaCompact}). The top of $\Db$ is the limit of its
own finite-stage approximants and is not $\omega$-compact
(\lean{BilimitGap.topSeq_not_omegaCompact}). An isomorphism sends top to top,
so $\BB\not\iso\Db$ (\lean{BilimitGap.no_orderIso}). And $\BB$ has a greatest
element while $\Dc$ has no maximal element (\cref{prop:nofootholds}), so
$\BB\not\iso\Dc$ (\lean{BilimitGap.no_orderIso_DInfty}).
\end{proof}

An isomorphism $\BB\iso\Db$ would have transported Theorem~B to $\BB$
(\lean{BilimitGap.rigid_of_orderIso}); \cref{thm:gap} closes that route.

\section{Theorem B: rigidity of \texorpdfstring{$\Dc$}{D-infinity} and \texorpdfstring{$\Db$}{its bilimit}}\label{sec:B}

\begin{theorem}\label{thm:B}
$\Aut(\Dc,\le)=\{\id\}$ and $\Aut(\Db,\le)=\{\id\}$.
\end{theorem}

\begin{proof}[Lean]
\lean{CopyCHigh.colimit_compat}, \lean{CopyCHigh.DBilimit_rigid}.
\end{proof}

The bilimit statement follows from the colimit one through \cref{thm:compacts}
(\lean{SierpinskiApplication.DBilimit_rigid_of_DInfty_rigid}). An automorphism
$g$ of $\Db$ preserves compactness, so it maps the image of $\Dc$ onto itself;
its restriction is an automorphism of $\Dc$, hence the identity. Then $g$ fixes
every compact element, and every element is the directed supremum of the compact
elements below it, so $g=\id$. The rest of this section sketches the rigidity of
$\Dc$.

\subsection{Finite levels}

\begin{theorem}\label{thm:finite}
Every $D_k$ is rigid. Consequently every monotone injection $D_k\to D_k$ is the
identity.
\end{theorem}

\begin{proof}[Sketch]
Induction on $k$ (\lean{SierpinskiApplication.D_all_rigid}). In the step, the
copy partition of $D_{k+1}$ is order-definable, so an automorphism restricts to
each copy, and each copy is isomorphic to $\End(D_k)$; it remains to show
$\End(P)$ is rigid when $P$ is. A finite-lattice automorphism is determined by
its action on join-irreducibles \cite{DaveyPriestley2002}, and the
join-irreducibles of $\End(P)$ form a poset isomorphic to
$P^{\mathrm{op}}\times\J(P)$ (\lean{GluingTowerLibrary.jiOrderIso}). Every
automorphism of this product splits as a product of automorphisms of the
factors (\cref{lem:prodaut}), each of which is the identity by rigidity of $P$
(\lean{EndPRigid.endP_rigid}). For a finite poset, rigidity implies that a
monotone injection, being a permutation with a monotone inverse power, is an
automorphism (\lean{MonotoneInjectionRigid.D_monotone_injection_eq_id}).
\end{proof}

\begin{lemma}\label{lem:prodaut}
Let $X,Y$ be finite atomic posets with bottom, such that $X$ has a unique atom
and $0<|Y|<|X|$. Then every automorphism $\varphi$ of $X\times Y$ splits:
$\varphi(a,b)=\bigl(\varphi(a,\bot)_1,\ \varphi(\bot,b)_2\bigr)$.
\end{lemma}

\begin{proof}[Comment]
This is \lean{ProdAut.prod_decomposition_of_card}. The proof counts atoms and up-sets: an automorphism cannot
exchange axes, because the factors are separated by cardinality and $X$ has a
unique atom. In the application $X=P^{\mathrm{op}}$ and $Y=\J(P)$. This is where
one might expect to invoke Sabidussi--Vizing unique factorization
\cite{Sabidussi1960,Vizing1963}; the order-theoretic setting makes it
unnecessary.
\end{proof}

\subsection{The colimit step}

The obstacle is \cref{prop:nofootholds}: $\Dc$ has no atom and no top, so no
invariant anchored at either end can tell which level a point comes from. The
argument proceeds in five steps.

\begin{enumerate}[label=\textbf{\arabic*.},leftmargin=2em]
\item \emph{Reduction to two points.} If every automorphism of $\Dc$ fixes the
  seam point $\mu_\infty$ and the image of the top of $D_0$, then it fixes the
  whole image of $D_0$ (\lean{TwoPointHypothesis}); by re-indexing (step~4) the
  same applies at every base level, which gives the theorem. This fixed-point claim is equivalent to a
  definability claim: the two points admit an order-theoretic description
  exactly when every automorphism fixes them (\lean{Definability}).
\item \emph{Exact shape of principal ideals.} The principal ideal of each
  join-irreducible point has an explicit product shape
  (\lean{SeedShapeProof.seedShape}). Comparing two points' ideals becomes a
  question about products of pairwise non-isomorphic finite factors, which no
  longer mentions the tower.
\item \emph{Truncation, not cancellation.} The comparison is settled by deleting
  the large factors with an order-theoretic predicate rather than cancelling a
  common infinite tail (\lean{Truncation}). This forces exact finite equations
  on a competitor's ideal: its cardinality, its atom count and an additive
  height must all match. The height is a natural-number-valued height proved
  additive on products, $\mathrm{ht}(a,b)=\mathrm{ht}(a)+\mathrm{ht}(b)$
  (\lean{Height}), since Mathlib's $\mathbb N_\infty$-valued height has no such
  lemma.
\item \emph{Re-indexing.} A comparison at a higher level is the same
  comparison at a shifted base, so the whole family reduces to one statement
  $\mathrm{FirstSlice}(n)$ for every base level $n$ (\lean{FinishLine}).
\item \emph{Three cases by copy.} $\mathrm{FirstSlice}(n)$ splits by the copy
  of the competing element. For copy $B$, and for the part of copy $C$ not above
  the marked point, the competitor's ideal is a non-trivial cut, while a
  product has only trivial cuts (\lean{CutObstruction}, \lean{FirstSliceCut}).
  In the remaining case, copy $C$ above the marked point, the competitor is
  necessarily a constant map $\cst q$ (\lean{CopyCHigh.copyCHigh_shape}), and
  its ideal splits by tag into three parts of sizes $|D_n|a$, $|D_n|a-1$ and
  $|D_n|b-1$ (\lean{CopyCHigh.card_partA}, \lean{CopyCHigh.card_partB},
  \lean{CopyCHigh.card_partC}). So the total is $\equiv-2\pmod{|D_n|}$. A
  product shape $D_n\times S$ would make it divisible by $|D_n|$, forcing
  $|D_n|\mid2$; but $|D_n|\ge3$ (\lean{CopyCHigh.three_le_card_D},
  \lean{CopyCHigh.not_card_dvd}). This closes the theorem.
\end{enumerate}

\section{Theorem A: \texorpdfstring{$\BB$}{B} is a fixed point of the gluing}\label{sec:A}

\begin{definition}
Let $\SEnd(\BB)$ be the poset of Scott-continuous self-maps of $\BB$, and
$\cst{\mu_\BB}$ the constant map at the interior point. Set
$\Phi_\sigma(\BB):=\Glue\bigl(\SEnd(\BB),\cst{\mu_\BB}\bigr)$, i.e.\
\cref{def:glue} with Scott-continuous maps in place of monotone ones
(\lean{PhiBilimit.markedB}, \lean{PhiBilimit.phiB}).
\end{definition}

\begin{theorem}\label{thm:A}
$\Phi_\sigma(\BB)\iso\BB$.
\end{theorem}

\begin{proof}[Lean]
\lean{PhiSurjective.theoremA_orderIso}.
\end{proof}

\begin{proof}[Sketch]
The argument is limit--colimit coincidence \cite{SmythPlotkin1982} adapted to
the gluing.
\begin{enumerate}[leftmargin=2em]
\item By construction $e_{n+1}=\Phi(\text{conjugation by }e_n)$ for all $n$, and
  $\Phi$ is a functor on tripointed maps (\cref{def:phi}). So the stages of
  $\Phi_\sigma(\BB)$ obtained by applying $\Phi$ to the stages of $\BB$ form a
  cocone, and three of the four colimit-cocone conditions are functoriality
  (\lean{PhiCocone}).
\item The fourth condition is that the images $\Phi(\delta_n)$ of the
  deflations $\delta_n$ of $\BB$ have supremum the identity. It follows from two
  general facts: composition preserves directed suprema in each argument, for
  Scott-continuous maps (\lean{EndLocalContinuity.comp_isLUB}); and the three
  copy inclusions into the glued carrier preserve least upper bounds with no
  side condition (\lean{GenContinuityTransport}, \lean{PhiCoconeLUB}).
\item The cocone yields the isomorphism once the relevant directed families in
  $\Phi_\sigma(\BB)$ have suprema and the cocone's projections preserve them.
  Both follow from a single explicit construction of the supremum together with
  its projections (\lean{PhiSurjective.exists_lub_with_proj},
  \lean{PhiSurjective.phiSupExists}, \lean{PhiSurjective.phiProjContinuous}).
\end{enumerate}
\end{proof}

\begin{remark}[Why Scott-continuous maps]
Step~2 uses Scott continuity essentially: for a monotone self-map $h$ that is
not Scott-continuous, the supremum of the approximants $\delta_nh\delta_n$ can
lie strictly below $h$. Taking the function-space step over $\SEnd$ rather than
$\End$ is what makes the identity $\bigvee_n\Phi(\delta_n)=\id$ available. At
finite levels the two coincide.
\end{remark}

\begin{remark}[A twisted action]
A variant action applies $f$ on copy $A$ and a second map $\psi$ on copies $B$
and $C$. It agrees with $\Phi(f)$ exactly when $\psi$ and $f$ agree away from
$\bot$ (\lean{PhiTriIdentification.phiTwist_eq_phiFun_iff}). For injective
tripointed self-maps of $\End(D_k)$ the two always agree, because $\End(D_k)$
is rigid and so its only monotone injection is the identity
(\lean{PhiTriIdentification.phiTwist_eq_phiFun_of_rigid},
\lean{PhiTriTower.phiTower_unique}). So on the tower the action of $\Phi$ on
morphisms is forced.
\end{remark}

Theorem~A says nothing about topology, about the gasket, or about $\Dc$ and
$\Db$.

\section{An internal IFS on \texorpdfstring{$\BB$}{B}, and the gasket}\label{sec:ifs}

Let $s\colon\Phi_\sigma(\BB)\to\BB$ be the isomorphism of Theorem~A
(\lean{InternalIFSInstance.sec}). For $j\in\{A,B,C\}$ put
$\sigma_j(x):=s\bigl(\tg{\cst x}j\bigr)$: embed $x$ as a constant map, include it
into copy $j$, and return to $\BB$. Each $\sigma_j$ is injective
(\lean{InternalIFSInstance.ifsB}).

\begin{proposition}\label{prop:escape}
$\Fix(\sigma_A)$ is not a singleton: $\sigma_A$ fixes both $\bot$ and
$\mu_\BB$, and $\sigma_B(\bot)=\mu_\BB$.
\end{proposition}
\begin{proof}[Lean]
\lean{EscapeAnalysis.fix_sigmaA_not_singleton}, \lean{EscapeAnalysis.sigmaA_mid},
\lean{EscapeAnalysis.sigmaB_bot}.
\end{proof}

So this system cannot underlie a coding map with one fixed point per copy. A
second system $(\sigma'_A,\sigma'_B,\sigma'_C)$ on the same carrier uses a
\emph{strict} constant map in place of $\cst x$, which keeps $\bot$ at $\bot$; it
has the same injectivity properties and $\mu_\BB\notin\Fix(\sigma'_A)$
(\lean{SigmaRebuild.ifsS}, \lean{SigmaRebuild.strictConst},
\lean{SigmaRebuild.mid_not_mem_fix}).

\begin{theorem}
$\Fix(\sigma'_A)=\{\bot\}$.
\end{theorem}
\begin{proof}[Lean]
\lean{EscapeProof.fix_sigmaA}.
\end{proof}

\subsection{The gasket and its coding map}
$\GG\subseteq[0,1]^2$ is constructed as the attractor of three just-touching
similarities, not assumed (\lean{Gasket.gasketSet_eq}). The address coding map
$\kappa\colon\{0,1,2\}^{\mathbb N}\to\GG$ is continuous and surjective, is the
unique map intertwining the shifts with the three similarities, and is a
quotient map (\lean{Gasket.continuous_codeGamma},
\lean{Gasket.codeGamma_surjective}, \lean{Gasket.coding_unique},
\lean{Gasket.isQuotientMap_codeGamma}).

\begin{proposition}
$\kappa$ is not an open map: the image of a single copy is not open in $\GG$.
\end{proposition}
\begin{proof}[Lean]
\lean{Gasket.not_isOpen_imageA}, \lean{Gasket.notOpen}.
\end{proof}

This concerns $\kappa$ only; we do not claim that every coding map onto $\GG$
fails to be open. It is the reason the interface for coding constructions in
the development (\lean{CoreIFS.TheoremAData}) asks only for a quotient map.

\subsection{Addresses in \texorpdfstring{$\BB$}{B} and the comparison map}
Using the repaired system, the address map $\beta\colon\{0,1,2\}^{\mathbb N}\to
\BB$ sends an address to the supremum of the chain it generates. It is
injective and intertwines the shift with the $\sigma'_j$; its image is the
\emph{attractor} $\mathcal A\subseteq\BB$. The comparison map
$\rho\colon\mathcal A\to\GG$ is defined by $\rho(\beta(w))=\kappa(w)$
(\lean{CodeB.codeB}, \lean{CodeB.codeB_injective}, \lean{CodeB.rho}).

\begin{theorem}
$\rho$ is surjective, it is the unique map $\mathcal A\to\GG$ intertwining the
$\sigma'_j$ with the three similarities, and it is chain-continuous: along a
monotone sequence in $\mathcal A$ with supremum $c$, the images converge to
$\rho(c)$.
\end{theorem}
\begin{proof}[Lean]
\lean{CodeB.rho_surjective}, \lean{AddressConvergence.rho_chainContinuous}.
\end{proof}

\begin{theorem}\label{thm:rhoneg}
$\rho$ is not injective. No topology on $\mathcal A$ whose open sets are upward
closed (in particular the Scott topology) makes $\rho$ continuous.
\end{theorem}
\begin{proof}[Lean]
\lean{SeamTriangle.rho_not_injective},
\lean{RhoContinuity.no_continuous_rho_of_upwardClosed}.
\end{proof}

The reason for non-injectivity is structural. Two addresses that name the same
seam point of $\GG$ reach different points of $\BB$: the address system of $\BB$
is \emph{seamless}, its copies meeting at no point
(\lean{SeamTriangle.codeB_seamless}), whereas the copies of $\GG$ meet at three
seams. Chain continuity is the form of continuity that survives.

\section{Automorphisms of \texorpdfstring{$\BB$}{B}: what is proved}\label{sec:Brigid}

We do not know whether $\BB$ is rigid. \Cref{thm:gap} rules out transporting
Theorem~B. We prove the following.

\begin{itemize}[leftmargin=2em]
\item An automorphism of $\BB$ that fixes every stage image pointwise is the
  identity (\lean{StageRigidity.eq_refl_of_fixes_stages}), and every level of the
  $\Phi$-iterate tower is maximally rigid
  (\lean{StageRigidity.levels_maximally_rigid}).
\item The $\omega$-compact elements of $\BB$ are exactly the stage images
  (\lean{PhiStagePipe.omegaCompact_iff_stage}). So every automorphism maps each
  stage image into some stage image and fixes the stage copies of $\bot$ and
  $\top$ (\lean{PhiStagePipe.stages_union_preserved},
  \lean{PhiStagePipe.apply_insD_top}).
\item $\BB$ has no atom (\lean{AtomlessBilimit.no_isAtom}).
\end{itemize}

\emph{Conditional rigidity.} Each of the following implies $\Aut(\BB)=\{\id\}$:
that every automorphism maps each stage image into the same stage
(\lean{StageRigidity.B_rigid_of_stageDefinable}); the same for cofinally many
stages (\lean{StageCofinal.B_rigid_of_cofinal}); and two uniqueness statements
for the level-to-level embeddings
(\lean{StageEmbedding.B_rigid_of_levelEmbUnique},
\lean{EmbDImage.B_rigid_of_levelImageStableJI}).

\emph{Refutations.} Both uniqueness statements are false. The constant inclusion
patched to send $\top$ to $\top$ is a second monotone, order-reflecting,
join-preserving, bound-preserving level embedding
(\lean{PatchedEmbedding.patched}, \lean{PatchedEmbedding.not_levelEmbUnique},
\lean{PatchedEmbedding.not_levelImageStableJI}); it does not extend to an
automorphism of $\BB$ (\lean{StageCover.no_patched_automorphism}). Stage
definability and its cofinal form remain open.

\emph{The seam point.} The canonical seam point $\chi$ of $\BB$ satisfies
$\mu_\BB<\chi<\top_\BB$. Whether every automorphism fixes it is equivalent to
stability of its principal ideal under all automorphisms
(\lean{SeamCut.cSeamFixed_iff_ideal_stable}), and follows if no other point has
an isomorphic principal ideal (\lean{IdealInvariant.cSeamFixed_of_ideal_pins}).
This too is open.

\section{The formalization}\label{sec:formal}

\subsection{Scope and size}
The release consists of 116 Lean modules, 37{,}109 lines, built against Lean
\texttt{v4.32.2} and Mathlib \texttt{v4.32.2}; the toolchain and every
dependency are pinned. The general order theory (gluing, $\End$,
join-irreducibles, products) is in \lean{GluingTowerLibrary} (10{,}386 lines);
the concrete tower and the constant-tower limits are in
\lean{SierpinskiApplication} (4{,}858 lines). The import graph is deep and
narrow, with a longest chain of 54 imports ending at \lean{Dinfty.Verify}, which
is imported by nothing. Pillar~A reuses Pillar~B's machinery:
\lean{PhiIterateTower} imports \lean{SeedShapeProof}, so the two theorems are
independent as statements but not as code.

\subsection{Trust base}
\lean{Dinfty/Verify.lean} restates each result of this paper as an anonymous
\texttt{example} whose proof is the named declaration, so a stated type cannot
drift from the proved one; any mismatch is a build error. It then runs
\texttt{\#print axioms} 195 times. In a clean build the reports are:
\begin{center}
\begin{tabular}{lr}
\toprule
axioms reported & count\\
\midrule
\texttt{propext}, \texttt{Classical.choice}, \texttt{Quot.sound} & 177\\
\texttt{propext}, \texttt{Quot.sound} & 9\\
\texttt{propext} & 6\\
none & 3\\
\midrule
total & 195\\
\bottomrule
\end{tabular}
\end{center}
No report contains \texttt{sorryAx} or \texttt{Lean.ofReduceBool}; the
development uses neither \texttt{sorry} nor \texttt{native\_decide}, and
declares no axioms of its own.

\subsection{Checking it}
With \texttt{elan} installed, the repository is checked by
\begin{center}
\texttt{lake exe cache get \&\& lake build Dinfty.Verify}
\end{center}
which takes roughly 15--20 minutes on a warm Mathlib cache. The repository
README gives an acceptance script that checks the build outcome and the axiom
distribution above, and the repository's continuous integration runs it on
every commit.

\subsection{Remarks on the development}
\begin{itemize}[leftmargin=2em]
\item \emph{Decidability.} The $\End$ lattices are finite but their order is
  decided classically, so tower levels cannot be enumerated by \texttt{decide};
  statements about the tower are proved structurally, uniformly in $k$, rather
  than by evaluation.
\item \emph{Negative results as theorems.} Each refutation in
  \cref{sec:ifs,sec:Brigid} is a Lean theorem with a negated conclusion, and the
  first IFS of \cref{prop:escape} is retained alongside the repaired one.
\item \emph{Reference labels.} Every declaration carries a stable label
  \texttt{[D-NNNN]}; \texttt{REFERENCE\_INDEX.md} in the repository maps labels
  to declarations and files.
\end{itemize}

\section{Conclusion}

A three-copy gluing of function-space lattices, iterated from the three-element
chain, yields limits with sharply different behaviour depending on how the
levels are glued together. Along the constant embeddings the colimit and bilimit
are rigid; along the embeddings generated by the gluing functor itself the
bilimit is a fixed point of the Scott-continuous gluing. The two limits are not
isomorphic, and the fixed point carries an internal IFS whose comparison with the
Sierpi\'nski gasket is surjective and chain-continuous but no better.

Two questions are left open: whether $\BB$ is rigid, and whether the
stage-definability hypothesis of \cref{sec:Brigid} holds. A positive answer to
the latter would settle the former.

\appendix
\section{Lean correspondence}\label{app:lean}

Module names are relative to \lean{Dinfty}; declaration names are as in the
release at \url{\repourl}, tag \texttt{\repotag}, archived at
\href{https://doi.org/\repodoi}{doi:\repodoi}.

\begin{longtable}{>{\raggedright\arraybackslash}p{0.34\textwidth}>{\raggedright\arraybackslash}p{0.58\textwidth}}
\toprule
Paper & Lean\\
\midrule
\endhead
$\Glue(Q,m)$, its order & \lean{GluingTowerLibrary.GlueCarGen}, \lean{GluingTowerLibrary.leGlue}, \lean{GluingTowerLibrary.rawRule}\\
$\End(P)$, $\SEnd(P)$ & \lean{GluingTowerLibrary.EndP}, \lean{GluingTowerLibrary.ScottEndP}\\
$D_k$, $\mu_k$, $\cst{\mu_k}$ & \lean{SierpinskiApplication.D}, \lean{SierpinskiApplication.dMid}, \lean{SierpinskiApplication.markedPt}\\
$\iota_k$, $p_k$ & \lean{SierpinskiApplication.iotaD}, \lean{SierpinskiApplication.pD}\\
$\Dc$, $\Db$ & \lean{SierpinskiApplication.DInfty}, \lean{SierpinskiApplication.DBilimit}\\
$\Phi$ on tripointed maps & \lean{PhiTriFunctor.phiFun}\\
$e_0$, $q_0$, $\lambda$ & \lean{SierpinskiApplication.glueAA}, \lean{GlueAAEpPartner.qAA}, \lean{GlueAAJoin.liftT}\\
$(e_n,q_n)$, $\BB$ & \lean{PhiIterateTower.iter}, \lean{PhiIterateTower.phiTowerEP}, \lean{PhiBilimit.B}\\
$\mu_\BB$, $\top_\BB$, $\Phi_\sigma(\BB)$ & \lean{PhiBilimit.midB}, \lean{PhiBilimit.topB}, \lean{PhiBilimit.phiB}\\
\cref{thm:compacts} & \lean{SierpinskiApplication.mem_range_toDBilimit_iff_isCompactElement}\\
\cref{prop:nofootholds} & \lean{DInftyNoFootholds.no_footholds}\\
\cref{thm:gap} & \lean{BilimitGap.no_orderIso}, \lean{BilimitGap.no_orderIso_DInfty}\\
Theorem B (\cref{thm:B}) & \lean{CopyCHigh.colimit_compat}, \lean{CopyCHigh.DBilimit_rigid}\\
\cref{thm:finite} & \lean{SierpinskiApplication.D_all_rigid}, \lean{MonotoneInjectionRigid.D_monotone_injection_eq_id}\\
\cref{lem:prodaut} & \lean{ProdAut.prod_decomposition_of_card}\\
Theorem A (\cref{thm:A}) & \lean{PhiSurjective.theoremA_orderIso}\\
$\sigma_j$; \cref{prop:escape} & \lean{InternalIFSInstance.ifsB}; \lean{EscapeAnalysis.fix_sigmaA_not_singleton}\\
$\sigma'_j$; $\Fix(\sigma'_A)=\{\bot\}$ & \lean{SigmaRebuild.ifsS}; \lean{EscapeProof.fix_sigmaA}\\
$\GG$, $\kappa$ & \lean{Gasket.Gamma}, \lean{Gasket.codeGamma}\\
$\kappa$ not open & \lean{Gasket.notOpen}\\
$\beta$, $\mathcal A$, $\rho$ & \lean{CodeB.codeB}, \lean{CodeB.Attractor}, \lean{CodeB.rho}\\
\cref{thm:rhoneg} & \lean{SeamTriangle.rho_not_injective}, \lean{RhoContinuity.no_continuous_rho_of_upwardClosed}\\
\bottomrule
\end{longtable}

\bibliographystyle{amsplain}
\bibliography{refs}

\end{document}
